% We switch to portrait mode. This works as advertised.
\documentclass[a0b,portrait]{a0poster}
% You might find the 'draft' option to a0 poster useful if you have
% lots of graphics, because they can take some time to process and
% display. (\documentclass[a0,draft]{a0poster})
\usepackage{graphicx}
% Switch off page numbers on a poster, obviously, and section numbers too.
\pagestyle{empty}
\setcounter{secnumdepth}{0}

% The textpos package is necessary to position textblocks at arbitary 
% places on the page.
\usepackage[absolute]{textpos}

% Graphics to include graphics. Times is nice on posters, but you
% might want to switch it off and go for CMR fonts.
\usepackage{graphicx,wrapfig,times}
\graphicspath{{images/}{../images}}

% These colours are tried and tested for titles and headers. Don't
% over use color!
\usepackage{xspace}
\usepackage{color}
\definecolor{DarkBlue}{rgb}{0.1,0.1,0.5}
\definecolor{Black}{rgb}{0.0,0.0,0.0}
\definecolor{Red}{rgb}{0.6,0.2,0} %{0.694,0.305,0.09}
\definecolor{LightGray}{rgb}{0.3,0.3,0.3}

\definecolor{ItalColor}{rgb}{0,0,0} %{0.2,0.6,0.2}
\definecolor{BoldColor}{rgb}{0,0,0} %{0.8,0.3,0.1}
\definecolor{TextColor}{rgb}{0,0,0} %{0.1,0.1,0.5}
\definecolor{CodeColor}{rgb}{0,0,0} %{0.1,0.1,0.1}
\definecolor{HeadColor}{rgb}{0.1,0.1,0.5}

% see documentation for a0poster class for the size options here
\let\Textsize\normalsize
\def\Head#1{\noindent\hbox to \hsize{\hfil{\huge\color{HeadColor} #1}}\bigskip}
\def\AHead#1{\bigskip\bigskip\noindent{\Large\color{HeadColor} #1}\smallskip}
\def\LHead#1{\bigskip\bigskip\noindent{\huge\color{HeadColor} #1}\smallskip}
\def\Subhead#1{\noindent{\huge\color{HeadColor} #1}}
\def\Title#1{\noindent{\VeryHuge\color{DarkBlue} #1}}

\TPGrid[40mm,40mm]{15}{25}  % 3 - 1 - 7 - 1 - 3 Columns

% Mess with these as you like
\parindent=0pt
%\parindent=1cm
\parskip=0.5\baselineskip

% abbreviations
\newcommand{\ddd}{\,\mathrm{d}}
\newcommand{\up}{\vspace*{-1em}}
\newcommand{\down}{\vspace*{1em}}
\newcommand{\I}[1]{{\color{ItalColor}\textit{#1}}\color{TextColor}}
\newcommand{\B}[1]{{\color{BoldColor}\textbf{#1}}\color{TextColor}}
\newcommand{\T}{\texttt}

\newcommand{\F}[1]{\B{FIXME: #1}}

\usepackage{srcltx}
\usepackage{fancyvrb}
\usepackage{ifpdf}

\DefineVerbatimEnvironment{mycode}{Verbatim}
{
  label=Code Example,
  fontsize=\normalsize,
  frame=single,
  framerule=1pt,
  framesep=1em,
% numbers=left,
  gobble=2
}

\begin{document}

\begin{textblock}{11.2}(0,0)
\baselineskip=3\baselineskip 
\Title{SAGA: The Simple API for Grid Applications\\
\Huge~Simple, Sustainable and Scalable Application Development}
\end{textblock}

\begin{textblock}{11.2}(0,1.35)
\AHead{Ole Weidner, Shantenu Jha\\[0.2em]
\large~Center for Computation \& Technology, Louisiana State University, Baton Rouge, U.S.A.
}
\end{textblock}

\begin{textblock}{3.8}(11.9,-0.2)
  \includegraphics[width=150mm]{cctfullcolortower}
\end{textblock}

\begin{textblock}{15}(0,2.6)
\color{LightGray}\rule{\linewidth}{2pt}
\end{textblock}

\begin{textblock}{15}(0,3.2)
  \includegraphics[width=788mm]{bigpicture}
\end{textblock}

\begin{textblock}{7}(0,8.5)

%%%%%%%%%%%%%%%%%%%%%%%%%%%%%%%%%%%%%%%%%%%%%%%%%%%%%%%%%%%%%%%%%%%%%%%%%%%%%%%%
\LHead{What is SAGA ?}
\large

\textbf{\color{DarkBlue} OGF Standard: } The Simple API for Grid Applications
(SAGA) is an official OGF standard (GFD-R-P.90) that describes a high level,
object-oriented API that can be used to develop distributed applications for
\textit{Clouds}, \textit{Grids} and other distributed infrastructures.  SAGA
provides abstractions for the three pilars of large-scale distributed
computing: job handling, data handling and communication. It is the only
community standard that tries to provide a comprehensive set of interfaces for
all three pillars with the same \textit{look and feel} (common syntactic and
semantic conventions), which tremendously simplifies application development.

To cope with the constantly changing landscape in distributed computing, SAGA
also describes a well-defined API extension process. This allows for the
addition of new interface groups (\textit{packages}) without compromising the
integrity of the API. The following API packages are currently defined in
SAGA:

\begin{itemize} \item{Job definition, submission and management API}
\item{File transfer, (remote) read/write API} \item{Replica catalog API}
\item{Advert catalog API} \item{Streaming client and server API} \item{Remote
procedure calls API} \end{itemize}

\textbf{\color{DarkBlue} SAGA Implementations: } There are currently several
independent SAGA implementations under development. The three most prominent
ones are CC-IN2P3's jSAGA (Java), Vrije University's JavaSAGA (Java, Python)
and Louisiana State University's SAGA implementation in C++, Python. All three
SAGA implementations consists of three components: (1) the API layer, (2) a
thin runtime layer and (3) \textit{adaptors} that translate the API calls to a
specific distributed middleware service (see figure). This 

%%%%%%%%%%%%%%%%%%%%%%%%%%%%%%%%%%%%%%%%%%%%%%%%%%%%%%%%%%%%%%%%%%%%%%%%%%%%%%%%
\LHead{Why Use SAGA ?}

\textbf{\color{DarkBlue} Shorter Development Time:} SAGA provides a very
simple, POSIX-inspired programming model with a consistent look and feel
across different API packages. It is very easy to learn and understand and
it saves the developer a lot of time by superseding the need to learning and
understanding different system-specific syntactic and semantic conventions.


\textbf{\color{DarkBlue} Portability:}

\textbf{\color{DarkBlue} Scalability:}

\large
\end{textblock}

\begin{textblock}{7}(8,8.5)

%%%%%%%%%%%%%%%%%%%%%%%%%%%%%%%%%%%%%%%%%%%%%%%%%%%%%%%%%%%%%%%%%%%%%%%%%%%%%%%%
\LHead{How SAGA is Used}
\large

There are many scientific computing projects that are successfully using SAGA
or SAGA-based software on national and international (or both!) grid and cloud
infrastructure to cary out production-level science and simulations. Examples 
include:

\textbf{\color{DarkBlue} High-Throughput Bio...?}

\textbf{\color{DarkBlue} DARE/RFOLD Biology Gateways}

\textbf{\color{DarkBlue} Cactus??}

\textbf{\color{DarkBlue} Ganga / Diane Framework } This project is a good
example of how SAGA can be used to extend existing frameworks and tools: Ganga
and Diane are tools frequently used by the HEP community at CERN to run
master-worker style applications on  gLite-based European grid infrastructure.
Adding SAGA-based job submission and file transfer capabilities to these tools
enabled researchers to run their applications on non gLite-based (e.g.
Globus-based) grids without having to change any of their existing applications.

\LHead{Road Ahead, \& Opportunities} 
\large



\end{textblock}

\begin{textblock}{15}(0,24.5)
\color{LightGray}
\rule{\linewidth}{2pt}
\end{textblock}

\end{document}

